%--------------------------------------------------------------------------------------
\thispagestyle{empty}
\chapter{Supersymmetric WKB}
The Wentzel–Kramers–Brillouin (WKB) approximation is a well known method of finding
asymptotic series in powers of $\hbar$ for the energy eigenvalues and eigenfunctions
of Schrödinger's equation: for bound states in a generalized potential well $V(x)$
which obeys the WKB condition
$$\left|\frac{\hbar p'(x)}{p^2}\right| =
\left|\frac{\hbar m\pdv[]{V}{x}}{p^3}\right|<<1\text{, where}
\; p(x):=\sqrt{2m\left[E-V(x)\right]}$$
i.e. slowly varying potentials and high momentum, the first semi-classical
correction yields a quantization condition for the first order of the energy
eigenvalues
\begin{equation}
\int_{x_1}^{x_2}\dd y \sqrt{2m\left[E_n-V(y)\right]}=\left(n+\frac{1}{2}\right)
\hbar\pi,
\label{eq:WKB_quant}
\end{equation}
with $V(x_1) = V(x_2) = E_n$.
It will be shown that this condition can yield more precise results using a
supersymmetric approach called the SWKB approximation.

\textbf{Note} In this chapter the $\hbar$ and $m$ dependencies of the potentials
will be explicit, since they are necessary for the WKB approximation.
\section{Lowest order SWKB: unbroken Supersymmetry}
Consider two supersymmetric partner potentials
$$
\begin{dcases*}
V_1(x) = W^2(x)-\frac{\hbar}{\sqrt{2m}}W'(x)\\
V_2(x) = W^2(x)+\frac{\hbar}{\sqrt{2m}}W'(x)
\end{dcases*}
$$
so that, formally, $W(x) = O(\hbar^0)$ and
$\hbar W'(x) = O(\hbar)$: the condition (\ref{eq:WKB_quant}) for $V_1$
is
$$\int_{x_1}^{x_2}\dd y\sqrt{2m\left[E^{(1)}_n-W^2(y)+\frac{\hbar}{\sqrt{2m}}W'(y)
\right]}=\left(n+\frac{1}{2}\right)\hbar\pi;
$$
expanding the integrand in powers of $\hbar$ yields
$$\sqrt{2m\left(E_n^{(1)}-W^2(y)\right)}+\frac{\hbar}{2}\frac{W'(y)}{\sqrt{E-W^2(y)}}+
O(\hbar^2),$$
the condition $V_1(x_1)=V_1(x_2)=E^{(1)}_n$ must be truncated at $O(\hbar^0)$
and so becomes
$$W^2(a)=W^2(b)=E_n^{(1)}\text{, where }
x_1 = a + O(\hbar),\;x_2 = b + O(\hbar).$$
The condition (\ref{eq:WKB_quant}) becomes
$$\int_{a}^{b}\dd y \sqrt{2m\left(E_n^{(1)}-W^2(y)\right)}+\frac{\hbar}{2}
\int_{a}^{b}\dd y \frac{W'(y)}{\sqrt{E^{(1)}_n-W^2(y)}}=\left(n+\frac{1}{2}\right)
\hbar\pi,$$
the second integral is elementar and its value is
$$\eval{\arcsin\left(\frac{W(y)}{\sqrt{E^{(1)}_n}}\right)}_a^b.$$
Supposing Supersymmetry is unbroken, conditions (\ref{eq:lim_W}) imply
$$-W(a)=W(b)=\sqrt{E^{(1)}_n}$$
that is
$$\eval{\arcsin\left(\frac{W(y)}{\sqrt{E^{(1)}_n}}\right)}_a^b=\pi.$$
So, with the hypothesis of unbroken Supersymmetry, the first order SWKB
quantization condition on $V_1(x)$ is
\begin{equation}
\int_{a}^{b}\dd y \sqrt{2m\left(E^{(1)}_n-W^2(y)\right)} = n\hbar\pi.
\label{eq:SWKB_quant_V1}
\end{equation}
The calculations for $V_2(x)$ are exactly the same, save a negative sign before the
$O(\hbar)$ term: its SWKB quantization condition is
\begin{equation}
\int_{a}^{b}\dd y \sqrt{2m\left(E^{(2)}_n-W^2(y)\right)} = (n+1)\hbar\pi.
\label{eq:SWKB_quant_V2}
\end{equation}
The advantages of this method are immediately clear: since the integrand
is strictly positive $\forall y\in[a,b]$, the quantization condition of $V_1$
with $n=0$ implies that the extrema of integration must coincide, so
$$W(a)=-W(b)=-W(a)=0=\sqrt{E^{(1)}_0}:$$
the ground state eigenvalue is exactly known to be
$E^{(1)}_0=0$ just from the first $\hbar$ order; moreover, these approximations
for the eigenvalues of the partner potentials respect the fundamental property of 
supersymmetric partner potentials $E^{(1)}_{n+1}=E^{(2)}_n$.

\textbf{Note} The regular WKB approximation yields better results the
greater is $n$, but this supersymmetric method is also exact for the ground state
energy: given this other constraint, it is natural to think that the SWKB method
will be, in general, a better approximation in the case of unbroken Supersymmetry.

\section{Lowest order SWKB: broken Supersymmetry}
The derivation of the new quantization condition for broken supersymmetry
is similar to the one for unbroken supersymmetry, with the key difference being
that at the $O(\hbar^0)$ turning points the superpotential must have the same
sign, as discussed after the definition of unbroken Supersymmetry:
$$W(a) = W(b) = \sqrt{E_n^{(1)}}$$
and so the $O(\hbar)$ integral vanishes for both $V_1$ and $V_2$, yielding
$$\int_a^b \dd y \sqrt{2m\left(E_n^{(i)}-W^2(y)\right)} = \left(n+\frac{1}{2}
\right)\hbar\pi$$\\

\textbf{Note} The reason why the SWKB method yields better results than regular
WKB for both unbroken and broken supersymmetry lies in the fact that the
factorization of supersymmetric partner Hamiltonians uncovers a term in the
potential that is $O(\hbar^1)$, that can be taken out of the quantization
condition integral and absorbed in the constant of the right hand side of the
regular WKB condition. For further discussion on broken supersymmetry,
see \cite{dutt_ref_in_ref}, as cited in \cite{cooper_main}
\section{SWKB for Shape Invariant Potentials}
This method is especially powerful for shape invariant potentials: consider the
family of potentials (\ref{eq:SIP_family})
$$V_s(x;a_1) = V_1(x;a_s) + \sum_{k=1}^{s-1}R(a_k):$$
the SWKB condition (\ref{eq:SWKB_quant_V1}) for the Hamiltonian $\hat H_1$ has the form
$$\int_{a}^{b}\dd y \sqrt{2m\left(E_n^{(1)}-W^2(y,a_s)\right)} = 
\int_{a}^{b}\dd y \sqrt{2m\left(E_n^{(s)}-\sum_{k=1}^{s-1}R(a_k)-W^2(x;a_s)\right)} = n\hbar\pi.$$
If supersymmetry is unbroken, as argued before, the quantization condition is
exact for the ground state energy
$$E^{(1)}_0 = 0 \iff E^{(s)}_0 = \sum_{k=1}^{s-1}R(a_k):$$
since all the potentials are shape invariant, this implies that the ground state energy \textit{for each
Hamiltonian} $\hat H_s$ is exact, moreover, for the properties discussed in the
previous chapter
$$E^{(s+1)}_0 = E_s^{(1)}$$
and so \textbf{all} the energy eigenvalues of ${\hat H_1}$ are exactly determined by the SWKB
quantization condition: this is a remarkable result, as the class of shape
invariant potentials is bigger than the class of known analytically solvable
potentials.




\COMMENT{
%explicit vs implicit \hbar dependency in W(x) and V_i(x), MF said it is unclear
\textbf{Note} This derivation is founded on the assumption that $W(x) =
O(\hbar^0)$ and $\hbar W'(x)=O(\hbar)$: this condition seems to imply that
this quantization condition holds only for potentials that contain an $O(\hbar)$
term, but it is always possible to absorb the $\hbar$ dependency with a choice of
dimensional parameters;
for example the superpotential for the reflectionless potentials yields
$$W(x;B) = B\tanh(x) \longrightarrow V_1(x;B) = B^2 - B(B+\hbar)\cosh^{-2}(x):$$
a choice of parameter $B$ so that $B(B+\hbar)$ is independent of $\hbar$ yields
an $O(\hbar^0)$ potential, with a constant $O(\hbar)$ energy
offset, which can always be discarded.}
